\section{Model: workload, resources, and value}
\label{sec:support}
\label{sec:allocation}
\subsection{Experiential substitution and implemented states}
Let $q\in[0,1]$ denote decision-relevant experiential fidelity and let
$\sigma(q)\in[0,1]$ be an increasing substitution map. For a residence
project with experiential benefit $E(n)$ and residual value $X(n)$,
\begin{equation}
 B_{\rm phys}(n,q)=[1-\sigma(q)]E(n)+X(n).
 \label{eq:experiential}
\end{equation}
Perfect substitution removes the first term, not external facts, materials,
geographical protection, or strategic options. For a discrete project with
substitutable value $e_j$, residual value $x_j$, and incremental resource
cost valued at $c_j$, the net comparison is
\begin{equation}
 \Delta_j(q)=[1-\sigma(q)]e_j+x_j-c_j.\label{eq:premium}
\end{equation}
A positive reality premium need not cover its opportunity cost. The
continuous model below specializes to $\sigma(q)=q$. This is a statement
about a decision-maker's substitution, not about the metaphysical identity
of a physical event and its simulation.

A regular hierarchy with branching factor $b\geq2$ and finite integer
depth $d\geq0$ contains
\begin{equation}
 W_d=\sum_{j=0}^{d}b^j=\frac{b^{d+1}-1}{b-1}.\label{eq:worlds}
\end{equation}
For per-world service requirement $u_k>0$, efficiency $\eta_k>0$, and
capacity $K_k$, homogeneous support requires
\begin{equation}
 D_k(d)=u_kW_d/\eta_k\leq K_k\quad\hbox{for every service }k.
 \label{eq:resource}
\end{equation}
Processing, state storage, heat rejection, and material replacement are
separate constraints. Their units cannot be added without an explicit
conversion. Cumulative consumption and available stocks must refer to the
same horizon; instantaneous power is not a lifetime energy supply.

\begin{lemma}[Finite-depth support frontier]
\label{prop:depth}
Let $Q=\min_kK_k\eta_k/u_k$. If $Q\geq1$, the largest feasible depth is
\begin{equation}
 d_{\max}=\left\lfloor\log_b[1+(b-1)Q]-1\right\rfloor.
 \label{eq:depth-bound}
\end{equation}
If $Q<1$, even the root world is infeasible at the specified workload.
\end{lemma}
\begin{proof}
The intersection of the service constraints is $W_d\leq Q$. Substituting
the geometric sum, taking the increasing base-$b$ logarithm, and restricting
to non-negative integers gives the result.
\end{proof}

For fixed per-world workloads,
$D_k(d+1)/D_k(d)=b+1/W_d\to b$. With $C_d=cW_d$ and logarithmic
value, the increment is
$\ln\{[1+c(bW_d+1)]/[1+cW_d]\}\to\ln b$.
These limits divide numerator and denominator by $W_d$. Recursion
therefore multiplies implemented demand while producing only a bounded
increment in logarithmic value per additional level. A rule that generates
unbounded depth is not an implementation of infinitely many simultaneous
independent states. Reusing code also does not erase distinct stored
histories. Scheduling fewer active worlds can save update work while
leaving a growing persistent-memory bill.

On a specified horizon $[0,T]$, define
\begin{equation}
 \chi_T=\sup_{0\leq t\leq T}\max_k\frac{D_k(t)}{K_k(t)}.
 \label{eq:chi}
\end{equation}
Local physical sufficiency means $\chi_T\leq1$. For a ratio
$\chi_{k0}\exp[(g_{D,k}-g_{K,k})t]$ initially below one and a positive
growth-rate difference, its crossing time is
\begin{equation}
 T_k=\frac{\ln(1/\chi_{k0})}{g_{D,k}-g_{K,k}}.\label{eq:crossing}
\end{equation}
This follows by solving for a ratio of one. It separates a successful
finite interval from a permanently feasible implementation. The archived
examples and their first memory failure are retained in the supplement.
The new results do not assume those examples establish indefinite support.

\subsection{A common resource budget}
Let $z,x,y,w\geq0$ be local investment, resident infrastructure,
industry-only capacity, and observation-only capacity. With local
investment as numeraire,
\begin{equation}
 z+p_nx+p_hy+p_mw\leq B,\qquad p_j>0,
 \label{eq:budget}
\end{equation}
where $B>0$ unless the degenerate zero-budget case is expressly mentioned.
A support constraint and an allocation constraint answer different
questions. A physically sufficient local path can still lose to a more
valuable external use, while an attractive local allocation can fail a
memory or heat requirement. Their conjunction, rather than either
calculation alone, describes an implementable inward path.
